\honest{Belief as Attire}{Eliezer Yudkowsky}{2007}

In earlier posts I described three kinds of improper belief. Here is a fourth: belief worn as a badge of belonging, like a priest's vestments or a skullcap. Robin Hanson gave me the image, and I call it ``belief as attire.'' \nb{The post gives no way to tell a belief worn as a badge from a belief the group holds for good reasons, beyond the quip ``you can tell, because the Enemy talks about it.''}

My example: the men who flew planes into the World Trade Center ``undoubtedly'' saw themselves as heroes fighting alien monsters, ``a la the movie \textsc{Independence Day}.'' I cite nothing they wrote or said. The comparison I offer is a Hollywood film about an alien invasion.

Only an inexperienced nerd, I say, would say this in an Alabama bar. The American thing to say is that the terrorists ``hate our freedom'' and committed a ``cowardly act.'' \nb{The first phrase echoes President Bush's speech to Congress nine days after the attacks: ``They hate our freedoms.''} In America, I say, you cannot call a suicide bomber's act heroic self-sacrifice. \nb{Within two weeks of the attacks, Bill Maher said on television that staying in the airplane was not cowardly, and Susan Sontag wrote in \textsc{The New Yorker} that ``they were not cowards.'' Both were attacked for it. So it could be said, at a price.}

Next I explain why people can be passionate about improper beliefs. Religious people who no longer really expect their religion to be true work hard to convince themselves that they are passionate, ``But it's not the same fire they had as a child.'' This is my ``impression,'' and ``I confess I'm not an expert here.'' \nb{The post gives no way to tell desperation from passion, so any passionate believer can be described as desperate.}

Belonging to a tribe, on the other hand, is easy and strong, like cheering for a sports team. ``People will die for it,'' I add. The beliefs a tribe wears are then spoken with the full passion of belonging.

In a footnote I call Republicans versus Democrats a ``swindle'' and a false dilemma, ``a topic for another time.''

\nb{The post applies the idea to Americans in an Alabama bar, the Enemy, religious people and supporters of both parties. It does not ask whether any belief of its author or its readers is worn this way.}
